\section{Data Preparation and Feature Engineering}
\label{sec:data}

\subsection{From nine tables to one row per order, split three ways}

The preparation pipeline is inherited unchanged from Phase 1
(\code{src/data\_load.py}). Three of the nine source tables --- items, payments
and reviews --- are one-to-many against orders, so each is aggregated to the
order grain \emph{before} joining; a direct merge yields 118{,}437 rows rather
than 99{,}441 and silently reweights every mean towards large baskets. The grain
is asserted after the join, and \code{customer\_unique\_id} is carried throughout
because \code{customer\_id} is regenerated per order and is not a person.
Restricting to the fully-operational window leaves \textbf{99{,}091} orders.
Eleven cleaning decisions with their evidence are recorded as
\code{CLEANING\_DECISIONS} in the loader and tabulated in the Phase 1 appendix;
\code{src/report\_stats.py} still reproduces every Phase 1 number byte-for-byte
after the Phase 2 additions, which \code{tests/test\_leakage.py} asserts.

Each period has exactly one job. \textbf{Train} (2017-01-01\,--\,2018-04-30,
73{,}008 orders) fits every model, with expanding-window cross-validation inside
it tuning hyperparameters. \textbf{Validation} (2018-05-01\,--\,2018-06-30,
13{,}156) sets every decision threshold and fits the margin-to-risk calibration.
\textbf{Test} (2018-07-01\,--\,2018-08-31, 12{,}927) is scored once, at the end,
for models already chosen. The split is chronological rather than
stratified-random because the target drifts, and additionally disjoint on the
customer: 2{,}997 people ordered more than once, so a purely temporal cut can
still put one of a person's orders on each side, and the \textbf{227} straddling
buyers are removed from a clean test subset used in
Section~\ref{sec:results-review}.

\subsection{Features, and the moment each becomes knowable}

Phase 1's 25 features carry forward unchanged: 6 calendar (month, weekday, hour,
weekend, business hours, the Black Friday window), 11 basket-and-money
(\code{log\_price}, \code{freight\_ratio}, \code{log\_weight}, volume, item and
seller counts, instalments), 4 geographic (\code{log\_distance},
\code{same\_state}, customer state, seller region), 2 for the promise and 2
catalogue. The design carries no exact duplicates --- Phase 1's first draft held
\code{total\_price} alongside \code{log\_price} and \code{total\_payment}, which
makes a linear baseline meaningless --- and \code{purchase\_year} is absent
because the split is chronological, so a linear term on it would extrapolate into
an unseen year. Phase 2 adds \textbf{eight track-record features}, taking the
checkout set to 33 and the at-delivery set to 36: a seller's mean lead time, late
rate, low-review rate and order count; a route's mean lead time; a category's
low-review rate; and the platform's own lead time and late rate. The at-delivery
regime additionally sees three columns describing how delivery actually went.

\paragraph{Why these features.}
Phase 1 chose the feature groups from the structure of the data; the Milestone 1
feedback asked for the literature behind them. \emph{Geography:}
origin--destination structure and distance are the backbone of published
delivery-time models, from JD.com's arrival forecasts \citep{salari2022delivery}
to last-mile travel-time predictors \citep{liu2021lastmile}. \emph{The promise:}
the quoted date is itself a decision that encodes what the platform's routing
knew \citep{cui2024promise, salari2022delivery}, which is why it is admissible
and useful at checkout. \emph{Basket, money and calendar:} order size, value and
freight are the order-level controls through which \citet{deshpande2023logistics}
identify the logistics--rating link, and demand peaks stress fulfilment
\citep{cui2024promise}, which Section~\ref{sec:problem} measures for Black
Friday. \emph{Track records:} the past performance of the parties to a
transaction is the standard predictor in supply-chain delay models
\citep{baryannis2019risk}, and an as-of aggregate of a key's earlier outcomes is
the time-respecting form of target encoding \citep{miccibarreca2001target}; the
general vocabulary of counts, ratios and calendar encodings follows
\citet{kuhn2019feature} and \citet{heaton2016feature}. The empirical screen was
Phase 1's correlation matrix (its Figure~8): lead time correlates 0.39 with
distance, 0.38 with the promised window, 0.17 with freight and 0.07 with weight;
promise and distance correlate 0.52 with each other, so both are kept and the
trees resolve the overlap; and the review score correlates $-0.33$ with lead
time and essentially zero with price or basket size -- the delivery, not the
purchase, drives the rating.

\subsection{Track records, and the trap inside them}
\label{sec:history}

A seller's history is the strongest thing a marketplace knows about an order at
checkout, and the easiest place here to leak, because of a distinction that is
easy to miss: \emph{past} means the earlier order had \textbf{closed} before this
one was placed --- not merely that it was placed earlier. A parcel bought last
week that has not yet arrived has no lead time yet, so ordering the history by
\emph{purchase} time would let every order average in outcomes that had not
happened. That is Kaufman et al.'s \emph{illegitimate feature}
\citep{kaufman2012leakage} in a convincing disguise, and it would be invisible in
the metrics because it would simply look like a better model. Every aggregate is
therefore an as-of join against the moment the earlier outcome became observable
--- \code{order\_delivered\_customer\_date} for delivery,
\code{review\_creation\_date} for reviews --- with a strict boundary.

Two further guards. A mean over one previous order is noise rather than a rate,
so below a floor of five prior orders the value is undefined and a fallback chain
takes over (seller $\rightarrow$ route $\rightarrow$ seller state $\rightarrow$
platform), after which 0.95\% of orders lack any history value. And cold start is
information rather than error: \code{hist\_seller\_n} carries how many past
orders the estimate rests on, so a model can discount a seller it has barely
seen. The platform features additionally use a \textbf{28-day trailing window}
rather than an expanding mean --- lead time halves across the analysis window, so
an average over all history reports roughly 13 days while the current month runs
at 8 --- and their range, 5.2 to 17.1 days, is a direct measure of the drift the
models must survive.

Because a leak here would be invisible in the metrics, the construction is
checked rather than reviewed: \code{tests/test\_leakage.py} rebuilds
\code{hist\_seller\_lead\_days} from first principles for 200 random orders and
requires exact agreement, and a second check confirms the platform tracker
\emph{lags} rather than anticipates (Appendix~\ref{app:repro}).

\subsection{Leakage: one rule per prediction moment}

Leakage here is not one rule but one per \emph{prediction moment}, because what
is known depends on when the model runs. \emph{At checkout} nothing about
fulfilment or the customer's opinion exists, so delivery dates and everything
derived from them, the review score and its timestamps,
\code{order\_approved\_at}, \code{shipping\_limit\_date} and
\code{order\_status} are forbidden. \emph{At delivery} only the review and its
timestamps remain forbidden --- the delivery outcome is now history, and the
brief itself lists ``delivery performance'' among the features for that problem.
\code{assert\_no\_leakage(features, regime)} enforces the correct set before any
model is fitted, is tested in both directions, and raises on an unrecognised
regime name, because passing the wrong regime is itself the mistake it exists to
catch.

\code{order\_estimated\_delivery\_date} is admissible at checkout and matters:
the customer is shown that promise when the order is placed, so it precedes the
outcome and encodes whatever Olist's own routing knew. It is a useful feature and
--- Section~\ref{sec:results-reg} --- a poor prediction. The cost of getting the
regime wrong is measurable: adding one forbidden column, the realised delivery
time, lifts the late-delivery model's ROC-AUC from 0.589 to \textbf{1.000}. That
is not a result; it is the target restated.
